% WORKING REFERENCE -- NOT PART OF THE SUBMITTED DOCUMENT.
%
% Do NOT \input this into the thesis. It is an aide-memoire: the premises that govern every feature
% decision in this project, collected in one place so they do not have to be re-derived each time
% they matter. Where an entry restates something argued properly elsewhere, the thesis version is
% the one that gets submitted -- this is the short form, kept self-contained on purpose.
%
% Each entry ends in a Consequence line. An entry without one is a note, not a concept to keep in
% mind, and should be cut.
%
% Citation keys resolve against thesis/refs.bib, so any entry can be promoted into thesis prose
% without re-finding its source.
%
% STANDING CONSTRAINT: no exact figure from the Copenhagen General Population Study appears here,
% or anywhere in thesis/, until data access.

\section*{Concepts to keep in mind}

\subsection*{A. What each dataset is, and what it licenses}

\begin{description}

\item[ImageCAS-X is a clinically referred cohort, not a screening one.]
Every scan in it was ordered because someone had a reason to order it, and the descriptor table
records coronary artery disease in roughly half the 800 usable cases
\cite{bransby2026imagecasx,zeng2023imagecas}.
(The exact split recorded in \texttt{CLAUDE.md}, 388 with and 412 without, comes from the abandoned
\texttt{week1} branch and has not been re-verified here.)
\textit{Consequence:} it licenses building the pipeline, verifying the segmentation against
annotations, and measuring the reproducibility of each feature. It cannot license a
normative description of healthy coronary topology --- the population is wrong for that. The
annotations are why the bias is worth accepting: centerlines carrying segment labels, connectivity
that is exact rather than approximated by a distance threshold, and a dominance column recorded independently of the labels.

\item[CGPS supplies what ImageCAS-X cannot, with one unknown that matters.]
Participants are not selected for disease, the study re-scans the same people over time, and it
carries adjudicated outcomes \cite{fuchs2023cgps}.
\textit{Consequence:} research question~2 needs the repeat-scan subset to establish a measurement
floor, and question~4 needs the outcome-linked subset. \textbf{Whether those two subsets overlap,
and by how much, is unknown.} If the overlap is small, reproducibility and outcome association are
two disjoint studies rather than one, and the plan's week ordering changes. Ask Phillip before the
tracks merge.

\end{description}

\subsection*{B. The biology that governs interpretation}

\begin{description}

\item[Coronary artery disease reshapes the geometry it would be predicted from.]
Examining 136 hearts at autopsy, Glagov et al.\ found the artery enlarges outward as plaque grows,
so the lumen does not measurably narrow until the lesion occupies about 40\% of the vessel, after
which it falls away sharply \cite{glagov1987remodeling}.
\textit{Consequence:} any shape derived from a lumen is measured after that compensation has already
begun, so a cross-sectional association between geometry and disease is partly the disease writing
its own predictor. A reference distribution has to come from people free of disease rather than
free of symptoms. This is the entire reason the work is split across two cohorts rather than one.

\item[Plaque is focal, and it follows low wall shear stress.]
The whole tree is exposed to the same cholesterol and the same blood pressure, yet lesions form in
particular places and not others \cite{chatzizisis2007ess}.
Tracing the flow through five human coronary trees after death put the plaques almost exclusively on
the outer wall of the vessels leaving a bifurcation, with the divider between them left clean, and
along the inner wall of curved segments --- exactly where shear was lowest \cite{asakura1990flow}.
\textit{Consequence:} bifurcation angle and tortuosity are not arbitrary descriptors; each one sets
the size and the position of a region where plaque is expected. But shear is not a single axis from
safe to dangerous --- over an existing plaque, high shear is associated with loss of the stabilizing
fibrous tissue \cite{bajraktari2021highwss}. A feature built only to flag low-shear geometry
captures half the mechanism, and the write-up has to say so.

\end{description}

\subsection*{C. What the annotations bound}

\begin{description}

\item[A branch count measures the tracing protocol as much as the anatomy.]
Septal perforators, acute marginals and nodal arteries were not traced at all. A side branch was not
traced where its distal diameter fell below 1 mm, or where its connection to the parent was unclear.
Where a side branch itself bifurcated only the larger was traced, unless the two were equal or the
smaller exceeded 1.8 mm. Where a \emph{main} vessel bifurcated, main-versus-side was decided by
course rather than size: the branch continuing along the atrioventricular groove is the LCX even
when the obtuse marginal is larger. Recorded from the dataset paper's supplementary protocol in
\texttt{CLAUDE.md}.
\textit{Consequence:} every branch-level feature is a protocol measurement. Any prevalence statement
has to name the protocol, and any comparison against another cohort --- CGPS included, where the
labeling protocol will differ and probably will not annotate all branches --- must treat the
protocol difference as a confounder rather than assume comparability. Note that septal perforators
are left-sided, so this is not a right-side restriction.

\item[The right side is annotated more coarsely than the left, deliberately.]
Only \texttt{RCA}, \texttt{R-PDA} and \texttt{R-PLA} exist on the right, and the catch-all
\texttt{Other} label never appears there.
\textit{Consequence:} any left-versus-right comparison of branch count or branching pattern must
state this, or it reports an annotation protocol as an anatomical finding.

\item[Image quality confounds any branch-count statistic.]
Quality is graded 0 to 4; grade 0 is exactly the 200 excluded cases, so the 800 usable ones span
grades 1 to 4. A lower-quality scan plausibly yields fewer traceable distal branches.
\textit{Consequence:} report branch-count distributions stratified by image quality, or demonstrate
that the association is flat across grades before pooling them.

\item[The delivered centerlines are smoothed, and the human effort is not in them.]
They are not the hand-traced centerlines. They were regenerated from the manually corrected lumen
masks by skeletonization and Gaussian smoothing ($\sigma = 0.5$ mm over a 5-vertex window), because
the traced centerlines did not run through the center of the corrected lumen. Segment names were
carried across by nearest match and reviewed by the lead analyst.
\textit{Consequence:} that smoothing is baked into any tortuosity computed from the delivered
points, so a tortuosity index measured here is not directly comparable to one measured on
centerlines that were never smoothed, and the reproducibility study has to treat smoothing as part of the
measurement chain. The trustworthy human effort sits in the masks and the segment names, not in the
centerline geometry.

\end{description}

\subsection*{Not covered here}

The purely technical traps live in \texttt{CLAUDE.md} and belong beside the code rather than beside
the argument: RAS versus LPS coordinate frames and the mirrored overlay that looks entirely
plausible; per-case anisotropic voxel spacing, which rules out measuring anything in voxel space;
the thirteen left sides carrying two ostia, eleven of which are the absent-left-main variant rather
than corruption; and the three cases whose centerlines contain a cycle.
